\documentclass[10pt,a4paper]{amsart}
\usepackage[margin=19mm]{geometry}
\usepackage{amsmath,amssymb,mathtools}
\usepackage[scr=rsfs,cal=euler]{mathalfa}
\usepackage{xfrac}
\usepackage[unicode,hidelinks,bookmarks=false]{hyperref}
\newcommand{\ie}{i.e.\ }
\newcommand{\cf}{cf.\ }
\newcommand{\resp}{respectively\ }
\newcommand{\Subsection}{\subsection}
\providecommand{\nobreakdash}{\nobreak\hbox{-}}
\setlength{\emergencystretch}{2em}
\multlinegap0pt
\usepackage{bm}
\usepackage{bbm}
\usepackage{dsfont}
\usepackage{xspace}
\usepackage{cancel}
\usepackage{mathtools}
\usepackage[shortlabels]{enumitem}
\usepackage{amsthm}
\makeatletter
\@ifundefined{theorem}{\newtheorem{theorem}{Theorem}}{}
\@ifundefined{lemma}{\newtheorem{lemma}[theorem]{Lemma}}{}
\makeatother
\title[Bouvier's Conjecture and Dimension Sequences of Unique Facto]{Bouvier's Conjecture and Dimension Sequences of Unique Factorization Domains}
\author{Viet-Hoang Tran, Thieu N. Vo, Tan M. Nguyen}
\date{Source: arXiv:2609.01887v2 (CC BY 4.0)}
\begin{document}
\maketitle

\noindent\emph{Source and licence.} Bouvier's Conjecture and Dimension Sequences of Unique Factorization Domains. Version arXiv:2609.01887v2 (CC BY 4.0), distributed under the Creative Commons Attribution 4.0 International licence, as recorded in the arXiv licence metadata for this exact version and shown on the abstract page https://arxiv.org/abs/2609.01887v2. Full rights notice: licenses/arxiv-2609.01887-CC-BY-4.0.txt.\par\smallskip

\noindent\textit{Editorial scope.} Counted unit: the Theorem whose source label is \texttt{thm:relative} in the article (source0.tex lines 1481--1505, source label \texttt{thm:relative}), delivered together with the article's own complete original proof of it (source0.tex lines 1507--1879, one \texttt{proof} environment of 373 lines). No mathematics is written, repaired or re-derived: statement and proof are the article's own text line for line. Required context retained uncounted: lem:residue-extension (lines 223--237); thm:capacity (lines 268--276); thm:ordinary (lines 560--576); lem:relative-hit (lines 1273--1280); eq:relative-weight-bound (lines 1306--1311); lem:relative-horizontal (lines 1361--1383); lem:relative-saturation (lines 1449--1460). Every definition, display and label of those lines is retained unchanged. Omitted from this package and not redistributed: 1--222, 238--267, 277--559, 577--1272, 1281--1305, 1312--1360, 1384--1448, 1461--1480, 1506--1506, 1880--2369. Nothing inside a retained excerpt is omitted or altered: every retained body is delivered exactly as the article's lines are written, so no omission receipt of any kind accompanies this package. Citations: the retained excerpts carry no citation command at all, so no bibliography entry of the article is transcribed and no reference is redistributed. Reference adaptation: none was needed and none was made; every reference inside the delivered unit resolves inside it, so no unresolved reference or citation is produced. Printed slips kept literal: no printed word, symbol, hypothesis or number of the counted statement or of its proof was altered. Layout and class: the article's own class is \texttt{amsart} with packages including fontenc, lmodern, microtype, geometry, amsmath, amssymb, amsfonts, mathtools, amsthm, booktabs, enumitem, xcolor, hyperref, cleveref; the wrapper is the lane's pinned \texttt{amsart} editorial wrapper at 10pt on A4, which restates the theorem-like environments the retained text needs and adds the article's own macro definitions that the retained text needs (none), with extra wrapper packages (bm, bbm, dsfont, xspace, cancel, mathtools, enumitem). The delivered numbering is the unit's own. Assets: no retained span contains a figure environment, a graphics-inclusion command, a PDF-page inclusion, a table or a TikZ picture; no image file is copied into this package, the package declares no asset dependency and no image file is redistributed with it. Licence: version arXiv:2609.01887v2 of this article is distributed under the Creative Commons Attribution 4.0 International licence, as recorded in the arXiv licence metadata for this exact version and shown on the abstract page https://arxiv.org/abs/2609.01887v2. A full rights notice is delivered at licenses/arxiv-2609.01887-CC-BY-4.0.txt. Characters: every retained excerpt is pure ASCII, so the delivered engine, XeLaTeX with its default Latin Modern text font, reports no missing character.
\begin{lemma}[Residues under field extension]\label{lem:residue-extension}
Let \(K\subseteq E\) be fields, let \(W\) be a valuation ring of
\(E\), and put \(V=W\cap K\). Then
\begin{equation}\label{eq:residue-extension}
 \operatorname{trdeg}_{\kappa(V)}\kappa(W)
       \leq\operatorname{trdeg}_K E.
\end{equation}
In particular an algebraic field extension induces an algebraic residue
extension. If \(k\subseteq E\), the valuation is trivial on \(k\),
and some element of \(E\) has positive value, then
\begin{equation}\label{eq:positive-value-bound}
 1+\operatorname{trdeg}_k\kappa(W)
       \leq\operatorname{trdeg}_k E.
\end{equation}
\end{lemma}

\begin{theorem}[The prime-flag formula]\label{thm:capacity}
For every finite-dimensional domain \(A\) and every \(n\geq0\),
\begin{equation}\label{eq:capacity-formula}
 \dim A[n]=n+\max_{0=\mathfrak p_0\subsetneq\cdots\subsetneq\mathfrak p_r}
 \left\{r+\sum_{i=0}^{r-1}
        \min\bigl(n,\tau_A(\mathfrak p_i,\mathfrak p_{i+1})\bigr)\right\}.
\end{equation}
The flag consisting only of zero is allowed. Flags need not be saturated.
\end{theorem}

\begin{theorem}\label{thm:ordinary}
Let \(m\ge1\), let \(d_1,\ldots,d_m\) be nonnegative integers, and let
\(t_{j,i}\in\mathbb N_0\cup\{\infty\}\), \(1\le i\le d_j\).
There is a countable semilocal UFD \(C\supset\mathbb Q\), with maximal
ideals \(\mathfrak m_1,\ldots,\mathfrak m_m\), and prime elements
\(\pi_j\in\mathfrak m_j\setminus\bigcup_{k\ne j}\mathfrak m_k\), such that,
writing \(C_j=C_{\mathfrak m_j}\), one has
\begin{equation}\label{eq:ordinary-realization}
 \dim C_j[n]=n+d_j+1+\sum_{i=1}^{d_j}\min(n,t_{j,i})
            =1+\dim(C_j/\pi_jC_j)[n]\qquad(n\ge0).
\end{equation}
Each \(C_j[1/\pi_j]\) is a PID, and \(\operatorname{Frac}C\) has a
discrete valuation \(\nu_j\) dominating \(C_j\), with
\(\nu_j(\pi_j)=1\). The residue field of \(C_j\) is \(\mathbb Q\).
The common fraction field and every nonzero prime quotient field of
every \(C_j[1/\pi_j]\) are rationally Hilbertian.
\end{theorem}

\begin{lemma}[Finite Hilbert fibers]\label{lem:relative-hit}
Let \(K\) be rationally Hilbertian.  If an integral affine
\(K\)-algebra \(E\) is finite over a polynomial subring
\(K[h_1,\ldots,h_r]\), then there are \(c_i\in\mathbb Q\) such that
\(E/(h_i-\delta_i c_i)_i\) is a finite field extension of \(K\),
where the \(\delta_i\in K^\times\) may be specified in advance.
The tuple can avoid any prescribed nonzero polynomial conditions.
\end{lemma}

\begin{equation}\label{eq:relative-weight-bound}
f^d+\sum_{j=1}^d c_j(x,z)f^{d-j}=0,
 \qquad
 \sum_i\alpha_i/\ell_i<j
 \quad\text{for every monomial }z^\alpha\text{ in }c_j.
\end{equation}

\begin{lemma}[Uniform horizontal modification]\label{lem:relative-horizontal}
Let \(B\) be a PID containing \(\mathbb Q\), allowing a field,
and put \(F=\operatorname{Frac}B\).  Suppose \(E\) is a UFD
finite locally free over \(B[x,z_1,\ldots,z_k]\), \(k\ge1\),
and \(E/aE\) is a domain for each prime element \(a\in B\).
Let \(f\) be a prime of \(E\) not associated with a lower prime.
Choose the distinct \(\ell_i\) larger than the coordinate degrees
of the norm of \(f\), and satisfying
\eqref{eq:relative-weight-bound} for its characteristic polynomial.
Put \(g_i=z_i^{\ell_i}-x^{\ell_iM_i+1}-b_i\), \(b_i\in B\).
If
\[
 E_F/(f,g_1,\ldots,g_k)=K'
\]
is a finite field extension of \(F\), then
\[
 E'=E[T_1,\ldots,T_k]/(fT_i-g_i)_i
\]
is a UFD with the same fraction field as \(E\), finite locally
free over \(B[x,T]\).  Every \(E'/aE'\) is a nonzero domain
containing \(E/aE\), and \(E'/fE'\) is a domain, torsionfree
over \(B\), with generic fiber \(K'[T]\).
\end{lemma}

\begin{lemma}[Saturation and factoriality]\label{lem:relative-saturation}
Let \(R\) be a domain, \(0\ne\pi\in R\), and
\(r_1,\ldots,r_k\in R[T]\).  If their reductions form a
regular sequence modulo \(\pi\), then
\[
 ((r_1,\ldots,r_k):\pi)=(r_1,\ldots,r_k).
\]
If its localization at \(\pi\) is a nonzero domain, the quotient
is a domain with the displayed exact special fiber.  No
Noetherian hypothesis is needed.  Also, an atomic domain with a
prime \(\pi\) and factorial localization at \(\pi\) is a UFD.
\end{lemma}

\begin{theorem}\label{thm:relative}
Let \(C\) be a countable finite-dimensional local UFD containing
\(\mathbb Q\), let \(\pi\) be prime, and put
\[
 D=C/\pi C,\quad B=C[1/\pi],\quad
 F=\operatorname{Frac}C,\quad L=\operatorname{Frac}D.
\]
Assume that \(B\) is a PID, that \(F\) and every \(B/aB\)
for a prime \(a\in B\) are rationally Hilbertian, and that a
discrete valuation \(\nu\) dominates \(C\), with
\(\nu(\pi)=1\).  Assume also
\begin{equation}\label{eq:relative-lower-profile}
\dim C[n]=1+\dim D[n]\qquad(n\ge0).
\end{equation}
For every \(t\in\mathbb N_{>0}\cup\{\infty\}\), there is
a countable local UFD \(A\supset C\), with the same residue
field and \(\pi\) prime, such that \(A[1/\pi]\) is a PID and
\begin{equation}\label{eq:relative-profile}
\dim A[n]=\dim C[n]+\max\{1,\min(n,t)\}\quad(n\ge0).
\end{equation}
Its fraction field is a rational function field in \(t\)
variables over \(F\), with countably many variables when
\(t=\infty\).  In particular the lower rings supplied by
Theorem~\ref{thm:ordinary} satisfy the required hypotheses.
\end{theorem}

\begin{proof}[Construction and proof]
First let \(t<\infty\), put \(k=t-1\), and work inside the
fixed field \(F^*=F(x,z_1,\ldots,z_k)\).  Extend \(\nu\)
by the Gauss rule giving all new variables value one.  Also
extend the valuation of the DVR \(C_{(\pi)}\) by that rule,
calling it \(w\).  Then
\begin{equation}\label{eq:relative-gauss}
\kappa(w)=L(x/\pi,z_1/\pi,\ldots,z_k/\pi),
 \qquad w(\pi)=w(x)=w(z_i)=1.
\end{equation}
Its restriction is zero on \(C\setminus\pi C\).  Thus \(w\)
will be centered at the upper augmentation prime, whereas
\(\nu\) dominates the full local ring.

Maintain a full local stage \(S\subset F^*\), and an affine
generic model \(E\), with these invariants:
\begin{enumerate}[label=(\roman*)]
\item \(S\) is a local UFD containing \(C\) locally, with the
same residue field; all top generators are positive for both
valuations and \(\nu\) dominates \(S\).
\item \(E\) is a UFD finite locally free over
\(B[x_s,z_{s,1},\ldots,z_{s,k}]\); every lower prime fiber
\(E/aE\) is a domain; and \(S[1/\pi]\) is a localization of \(E\).
\item The special fiber is the local ring at the lower maximal
ideal and top origin of \(\Gamma_s[z_{s,1},\ldots,z_{s,k}]\).
Here \(\Gamma_s\) is a domain finite free over \(D[x_s]\),
augmented onto \(D\), its augmentation ideal is nilpotent
modulo \(x_s\), and
\[
 \operatorname{Frac}(L\Gamma_s)=L(u_s),\qquad
 x_s=\lambda_su_s^{d_s},\quad \lambda_s\in L^\times, d_s>0.
\]
\end{enumerate}
Initially these are the local polynomial ring \(C[x,z]\), the
affine ring \(B[x,z]\), and \(\Gamma_0=D[x]\).  Polynomial
rings over UFDs are factorial without a Noetherian hypothesis.
The generic models are Noetherian because \(B\) is a PID; the
full stages need not be Noetherian.

\smallskip\noindent\emph{Horizontal actions.}
For \(k>0\), choose a currently horizontal irreducible in
\(S[1/\pi]\), represented by a nonunit of \(S\).  The affine
UFD \(E\) supplies an associated affine prime \(f\), not
associated with a lower prime.  Its coefficients lie in
\(C[1/\pi]\), so multiply it by a power of \(\pi\) to put
it in \(S\), and divide out its largest \(\pi\)-power there.
It stays an affine prime of \(E\), belongs to \(S\), is not
\(\pi\)-divisible, and has positive \(\nu\)-value.

Choose distinct primes \(\ell_i\) satisfying all bounds in
Lemma~\ref{lem:relative-horizontal}, exceeding the degrees in the
passive variables of a numerator of \(\bar f\), and not dividing
\(d_s\).  Increase them, then choose \(N_i=\ell_iM_i+1\)
and integers \(q_i\) so that, for \(v=\nu,w\),
\begin{equation}\label{eq:relative-positive}
N_iv(x_s)>\ell_i v(z_{s,i}),\quad
 q_i>\ell_i v(z_{s,i})>v(\pi f).
\end{equation}
There are arbitrarily large permissible primes.  The norm
normalization proves that \(E_F/fE_F\) is finite over the
polynomial algebra in \(z_{s,i}^{\ell_i}-x_s^{N_i}\).
Lemma~\ref{lem:relative-hit} therefore chooses
\(c_i\in\mathbb Q^\times\) giving a finite field fiber, with
the prescribed factors \(\pi^{q_i}\).  Adjoin the actual fractions
\begin{equation}\label{eq:relative-horizontal-fractions}
Z_i=\frac{z_{s,i}^{\ell_i}-x_s^{N_i}-c_i\pi^{q_i}}{\pi f}
\end{equation}
and localize at the old maximal ideal and the new coordinates.
Lemma~\ref{lem:relative-horizontal}, with \(\pi f\) in place
of the associate \(f\), gives the new generic affine invariants.

We justify the full presentation, where \(C\) may be
non-Noetherian.  Its reduced equations are the successive monic
equations \(z_{s,i}^{\ell_i}=x_s^{N_i}\).  Each quotient is a
domain: over \(L(u_s)\), the exponent \(d_sN_i\) is coprime
to \(\ell_i\), and a valuation at \(u_s=0\) proves
irreducibility.  At later equations the old exponent is multiplied
by the preceding distinct primes, so the same coprimality holds.
Monic freeness gives the injection before extending coefficients
to \(L\).  In particular the equations are a regular sequence.
The bound on the passive degrees ensures that \(\bar f\) remains
nonzero: its numerator is already a nonzero standard monomial
combination for the successive monic equations, and its local
denominator still has unit augmentation.

Lemma~\ref{lem:relative-saturation} makes the full presentation
\(\pi\)-torsionfree.  Its generic localization is a localization
of the affine domain just proved.  The fractions
\eqref{eq:relative-horizontal-fractions} have positive value by
\eqref{eq:relative-positive}; every intended inverted denominator
therefore evaluates to a nonzero \(\nu\)-unit.  This proves the
localization is nonzero and the full presentation is its asserted
subring of \(F^*\), with exactly the displayed domain as special
fiber.  Its elements have bounded factorization length by the
integer \(\nu\)-value.  It is atomic, and the last part of
Lemma~\ref{lem:relative-saturation} makes it a UFD.

The new curve algebra remains finite free over \(D[x_s]\) and
its augmentation ideal is nilpotent modulo \(x_s\).  Rationality
of its generic field does not require roots of coefficient units:
if \(z^\ell=\lambda^N u^{dN}\) and
\(a\ell+b dN=1\), put \(v=u^a z^b\).  Then
\(u=\lambda^{-Nb}v^\ell\) and
\(z=\lambda^{Na}v^{dN}\).  Iteration supplies the asserted new
rational parameter.  The old generic quotient by \(f\), after
inverting the nonzero lower coefficients, has image in the finite
field chosen above.

\smallskip\noindent\emph{Actions at a lower prime.}
Let \(a\) be a prime of \(B\), represented by a prime of \(C\)
different from \(\pi\).  It is prime in \(E\) by (ii).
Choose distinct large prime exponents \(m_i\), coprime to the
current curve exponent, put \(N_i=m_iM_i+1\), fix \(e\ge1\),
and choose \(\ell\) and all \(q_i\) sufficiently large.  Set
\begin{equation}\label{eq:relative-mixed}
\begin{gathered}
g_0=x_s^\ell-c_0\pi^{q_0},\quad
 g_i=z_{s,i}^{m_i}-x_s^{N_i}-c_i\pi^{q_i},\\
 X=g_0/a^e,\qquad Z_i=g_i/(\pi a).
\end{gathered}
\end{equation}
Choose the exponents so that each first displayed monomial has
strictly smaller value than the other terms and strictly larger
value than its denominator, for both \(\nu\) and \(w\).
This is possible because \(e\) is fixed and the old coordinates
are positive.  The functions
\(x_s^\ell,z_{s,i}^{m_i}-x_s^{N_i}\) form a finite polynomial
normalization of \(E/aE\): the old \(x_s,z_{s,i}\) are integral
over them by their monic equations.  The Hilbert property of
\(B/aB\) chooses rational nonzero \(c_i\) with
\begin{equation}\label{eq:relative-mixed-field}
E/(a,g_0,g_1,\ldots,g_k)\text{ a finite field over }B/aB.
\end{equation}

For clarity, the affine presentation of this mixed action has no
hidden torsion.  At a prime containing \(a\), condition
\eqref{eq:relative-mixed-field} makes
\(a,g_0,\ldots,g_k\) a parameter sequence in the old
Cohen--Macaulay local ring.  In the polynomial local ring replace
the \(g_i\) by \(a^eX-g_0,\pi aZ_i-g_i\); after \(a\)
these are the same regular sequence up to sign.  Permuting it
makes \(a\) regular on their quotient.  Away from \(a\) the
equations eliminate the new coordinates.  The quotient thus
embeds in \(E[1/a]\), is a domain, and is Cohen--Macaulay.  Its
quotient by \(a\) is a polynomial ring over the field in
\eqref{eq:relative-mixed-field}, so \(a\) is prime.  Every other
lower prime \(b\) has \(a\) invertible in \(B/bB\), so the
new \(b\)-fiber is just \(E/bE\) with the fractions evaluated;
it is a domain.  Nagata gives a UFD.  The equations
\[
 x_s^\ell=a^eX+c_0\pi^{q_0},\qquad
 z_{s,i}^{m_i}=x_s^{N_i}+c_i\pi^{q_i}+\pi aZ_i
\]
make it finite over \(B[X,Z]\).  The unchanged fraction field
proves these new base coordinates independent.  The
Cohen--Macaulay and regular-base argument in
Lemma~\ref{lem:relative-horizontal} proves local freeness.

Its exact full special fiber also needs a non-Noetherian check.
First impose the monic equations on the old passive block, giving
a domain \(\Gamma_1\) free and finite over \(D[x_s]\).  The
remaining equation is \((\bar a)^eX=x_s^\ell\).  The quotient
\(\Gamma_1[X]/((\bar a)^eX-x_s^\ell)\) is
\(\bar a\)-torsionfree.  Indeed, reduce
\(\bar a H=((\bar a)^eX-x_s^\ell)G\) modulo \(\bar a\).
Multiplication by \(x_s^\ell\) is injective on
\(\Gamma_1/\bar a\Gamma_1\), since \(\Gamma_1\) is free
over \(D[x_s]\).  Thus \(G\) is \(\bar a\)-divisible,
and cancellation proves saturation.  Inverting \(\bar a\)
gives a domain, so the quotient is a domain.  It is finite free
over \(D[X]\), using the monic equation for \(x_s\).  All old
augmentation elements are nilpotent modulo \(X\), and its
generic field is the same rational curve with
\(X=x_s^\ell/(\bar a)^e\).

The full reduced sequence is the successive monic equations
followed by this nonzero equation over a domain.  Apply
Lemma~\ref{lem:relative-saturation} to obtain \(\pi\)-saturation.
The actual positive fractions in \eqref{eq:relative-mixed} ensure
nonzero localization.  Atomicity and the generic UFD then prove
that the full stage is a UFD and \(\pi\) is prime.  A lower
prime cannot be lost on localizing: its value under \(\nu\) is
positive; if it became a unit after inverting \(\pi\), the UFD
would make it a \(\pi\)-power times a unit, contrary to its
nonzero special reduction.  The same reasoning retains the
chosen horizontal prime.  This completes preservation of all
stage invariants, including local survival of the affine fibers.

\smallskip\noindent\emph{Fair scheduling and factoriality of the union.}
Enumerate the elements of all countably many full stages and
repeat their tasks.  When an element is currently an irreducible
horizontal nonunit, perform its horizontal action; perform a
mixed action at each lower prime cofinally.  Each final
irreducible is irreducible in every later full stage containing
it, since nonunits remain nonunits under the dominating valuation.
Thus an eventual horizontal irreducible receives arbitrarily
late actions.  The integer valuation bounds every factorization
in the union \(A\), so it is atomic; an irreducible is prime
in all later full UFD stages, hence is prime in the union.
Therefore \(A\) is a local UFD.  The domain special fibers
and the preserved lower primes remain domains at the limit;
noninjective special-fiber transition maps cause no difficulty,
since products are witnessed at a finite stage.

Write \(G=A[1/\pi]\).  Repeated mixed actions put every old
element of \(G/aG\) into a finite field over \(B/aB\).
An old denominator has nonzero image in this field, since otherwise
the localized quotient would be zero and \(a\) would be a unit.
Thus this includes old localized elements.
Its nonzero elements keep their inverses under subsequent maps,
so \(G/aG\) is algebraic over that field and hence is itself
a field.  After inverting \(B\setminus\{0\}\), repeated
horizontal actions have the same consequence for every
horizontal prime quotient, now over \(F\).  Thus the latter
localization has dimension at most one.  If \(k=0\), this
conclusion requires no horizontal action: every nonzero prime
quotient of each affine generic curve over \(F\) is already
algebraic over \(F\), and a persistent nonzero prime witness
retains that property in every later chart.

A nonzero horizontal prime cannot lie below \(aG\), since
the latter is principal of height one in a UFD.  These facts
give \(\dim G\le1\).  A UFD of dimension at most one is a
PID: the gcd of a nonzero ideal is the gcd of finitely many
of its elements, since one fixed nonzero element has only
finitely many prime factors.  Divide the ideal by that gcd.
If the resulting ideal were proper, a maximal ideal above it
would be generated by a prime element (every nonzero prime
has height one), contradicting that the finite gcd is one.
Consequently a height-one prime \(q\) of \(A\) avoiding \(\pi\)
has quotient fraction field algebraic over \(F\) if \(q\cap C=0\),
or algebraic over \(B/aB=\operatorname{Frac}(C/aC)\) if its
contraction is the lower prime \((a)\).

\smallskip\noindent\emph{The boundary.}
Let \(\mathcal D=A/\pi A\), let
\(\epsilon:\mathcal D\to D\) be the augmentation, and put
\(M=\ker\epsilon\).  Every passive block is eventually
processed into a later curve.  To justify arbitrarily late
actions when \(k>0\), the nonzero old curve coordinate has
nonzero special image in every later curve and has positive
\(\nu\)-value.  A final irreducible factor other than \(\pi\)
is either horizontal or a lower prime; its repeated queue
therefore supplies such actions.  For \(k=0\) there is no
passive block to consume.

Every finite subset of \(L\mathcal D\) consequently lies in a
localization of an affine curve with rational fraction field
over \(L\).  Finite witnesses of strict prime inclusions show
that this ring has dimension at most one.  Its nonzero
augmentation prime has residue field \(L\), giving equality.
A persistent nonzero prime witness in a curve chart makes all
later residue fields algebraic over \(L\); hence every nonzero
prime of \(L\mathcal D\) has algebraic residue field.  Its
fraction field has transcendence degree one over \(L\): old
curve fields inject into the later curve fields, each finitely
many passive coordinates eventually becomes algebraic over an
old curve parameter, and that parameter stays transcendental.

If \(Q\) is a prime of \(\mathcal D\) contracting to
\(0\ne p\subset D\), lift a nonzero element of \(p\) to
\(C\) and factor it there.  Some prime factor \(a\ne\pi\)
has nonzero reduction in \(p\).  At a mixed \(a\)-action the
equation \((\bar a)^eX=x_s^\ell\) forces \(x_s\) into
\(Q\); nilpotence modulo \(x_s\) forces all earlier curve
augmentation elements into \(Q\).  Repeating after every
passive block has been consumed gives \(M\subset Q\).  Thus
\begin{equation}\label{eq:relative-boundary-primes}
Q=\epsilon^{-1}(p)\quad\text{if }Q\cap D=p\ne0.
\end{equation}

Here are the dimension details needed from these boundary facts.
For a nonzero zero-contraction boundary prime \(P\), the first
capacity \(\tau_{\mathcal D}(0,P)\) is zero: a centered
valuation is nontrivial and trivial on \(L\), and one positive
element together with independent residue lifts are algebraically
independent over \(L\).  A field of transcendence degree one
therefore leaves no residue transcendence.  Also
\(\operatorname{Frac}(\mathcal D/P)/L\) is algebraic.  Write
\(P_p=\epsilon^{-1}(p)\) for \(p\ne0\).  Capacities between
two such primes are exactly those of \(D\), because the
quotients are \(D/p\).  Lemma~\ref{lem:residue-extension} gives
\[
 \tau_{\mathcal D}(0,P_p)\le1+\tau_D(0,p),\qquad
 \tau_{\mathcal D}(P,P_p)\le\tau_D(0,p).
\]
A flag has at most one nonzero zero-contraction prime.  If it
jumps directly to \(P_p\), use
\(1+\min(n,1+\tau)\le2+\min(n,\tau)\); if it first
passes through \(P\), use its zero first capacity and the
second inequality.  All later contributions are those of a
flag of \(D\).  Theorem~\ref{thm:capacity} gives the upper
bound below.  For the lower bound lift a chain through the
surjection to \(D[n]\) and prepend zero below \(M[n]\).
Consequently
\begin{equation}\label{eq:relative-boundary-profile}
\dim\mathcal D[n]=1+\dim D[n].
\end{equation}

\smallskip\noindent\emph{Capacity and exact profile.}
Let \(h\) be the inverse image of \(M\) in \(A\).  Its
contraction to \(C\) is \((\pi)\), with residue field \(L\).
The valuation \(w\) dominates \(A_h\): a denominator outside
\(h\) has nonzero lower augmentation, hence value zero.
Equation~\eqref{eq:relative-gauss} gives \(\tau_A(0,h)\ge t\).
Conversely a valuation centered there restricts on \(F\) to
the DVR \(C_{(\pi)}\).  The residue extension inequality and
\(\operatorname{trdeg}_F F^*=t\) give the reverse inequality.
Thus \(\tau_A(0,h)=t\).  The prime \(\pi A\) has height
one, the generic ring is a PID, and the boundary localization
at \(M\) has dimension one, so \(\operatorname{ht}h=2\).

The prime \(\pi\) and the entire boundary give
\(\dim A[n]\ge\dim C[n]+1\).  Concatenating the edge
\(0<h\), of capacity \(t\), with a flag of \(A/h=D\)
gives \(\dim A[n]\ge\dim C[n]+\min(n,t)\).
We prove the reverse bound explicitly.  The prime classification
above makes \(A\) finite-dimensional: a flag contains at most
one nonzero prime avoiding \(\pi\), followed by a boundary
flag.  Apply Theorem~\ref{thm:capacity} to flags beginning at zero.

If the first nonzero prime is \(\pi A\), its first capacity
is zero since \(A_{(\pi)}\) is a DVR; the rest is bounded by
\eqref{eq:relative-boundary-profile}, giving \(\dim C[n]+1\).
If that first prime is \(q\) avoiding \(\pi\), it has height
one and first capacity zero.  For \(q\cap C=0\), delete this
one initial edge and contract the rest to \(C\).  The
contractions are strict: at most one nonzero generic-curve
prime lies over \((\pi)\), and higher primes are the distinct
augmentation preimages in \eqref{eq:relative-boundary-primes}.
The field \(\operatorname{Frac}(A/q)\) is algebraic over
\(F\); subsequent fields are lower quotient fields or
algebraic extensions of them.  The residue extension inequality
bounds every remaining capacity by its lower counterpart.  This
case contributes at most \(\dim C[n]+1\).  If instead
\(q\cap C=(a)\ne0\), all subsequent lower contractions
contain both \(a\) and \(\pi\), so are augmentation preimages.
Contract the whole flag to \(0<(a)<\cdots\); algebraicity of
the quotient field over \(\operatorname{Frac}(C/aC)\) bounds
all capacities without any extra edge.  The bound is \(\dim C[n]\).

Finally suppose the first nonzero prime \(Q\) properly contains
\(\pi\).  Its contraction \(c\) is nonzero, and contractions
of the whole flag are strict for the same boundary reason.
Only the first edge can gain capacity: restricting a centered
valuation from \(F^*\) to \(F\) gives
\[
 \tau_A(0,Q)\le t+\tau_C(0,c),\qquad
 \min(n,\tau_A(0,Q))
 \le\min(n,t)+\min(n,\tau_C(0,c)).
\]
The endpoint residue extension is algebraic for a generic-curve
prime, and identical for an augmentation preimage.  Every later
capacity is bounded by the corresponding contracted one.  The
flag lengths agree, giving the bound \(\dim C[n]+\min(n,t)\).
These cases prove \eqref{eq:relative-profile}.

\smallskip\noindent\emph{Infinite capacity.}
For \(t=\infty\), start with \(F(x,z_1,z_2,\ldots)\) and
the two Gauss valuations, still integer-valued.  Introduce one
fresh independent positive variable at each round; each stage
uses a finite tuple.  At every horizontal or mixed action process
the entire current tuple, with exponents chosen for that finite
stage.  The finite affine and saturation proofs apply literally.
Each old variable is consumed into a later boundary curve, and
each old generic quotient element is covered by a later action
at its final prime.  Thus the same boundary, factorization, and
generic PID arguments hold.  The fixed valuation \(w\) now has
residue transcendence degree infinity over \(L\), so
\(\tau_A(0,h)=\infty\).  In the last flag bound simply use
\(\min(n,\tau)\le n\); all other cases are unchanged.  This
proves the theorem for infinite capacity as well.
\end{proof}

\end{document}
